% Copy-paste in Paper (Pfad B, nach QQ-Spektraltabelle)
\begin{definition}[Arithmetischer Schalen-Kollaps über Hurwitz]
Sei $H$ die Hurwitz-Maximalordnung in $\mathbb{H}=\mathrm{QuaternionAlgebra}(\mathbb{Q},-1,-1)$
mit $E,A,B,C \leftrightarrow 1,i,j,k$. Für den Rundweg-Operator
$\Pi_\Gamma=\frac14(I+\Gamma+\Gamma^2+\Gamma^3)$ auf $\mathbb{Q}^4$ sei
$V_{\lambda=0}=\ker(\Pi_\Gamma)$ der antisymmetrische Schalen-Träger
(z.\,B.\ $q_{\mathrm{asym}}=i-k$).
Für eine Primquaternion $\pi_p\in H$ mit $N(\pi_p)=p$ und das Linksideal $I_p=H\cdot\pi_p$
heißt $q\in H$ ein \emph{arithmetischer Resonanzpunkt mod $p$}, wenn
$0\neq q\in V_{\lambda=0}\cap I_p$.
\end{definition}
\begin{proposition}[Operative Kriterien]
Mit $q=h\pi_p$ gilt $q\in I_p$ genau dann, wenn $h=q\pi_p^{-1}$ Hurwitz ist;
notwendig ist $p\mid N(q)$. Für $q_{\mathrm{asym}}=i-k$ und $p\in\{5,7\}$
mit $\pi_5=2+i$, $\pi_7=2+i+j+k$ folgt aus der Sage-Verifikation
\texttt{rundweg\_hurwitz\_resonanz.sage}, dass kein nichttriviales rationales
Vielfach von $q_{\mathrm{asym}}$ in $I_5\cup I_7$ liegt.
\end{proposition}
